\documentclass[11pt]{article}
\usepackage[margin=0.85in]{geometry}
\usepackage{amsmath,amssymb,amsthm}
\usepackage{booktabs}
\usepackage{longtable}
\usepackage{array}
\usepackage[T1]{fontenc}
\usepackage{lmodern}

\newcommand{\Fix}{\operatorname{Fix}}
\newcommand{\Stab}{\operatorname{Stab}}
\newcommand{\Pow}{\mathcal{P}}
\newcommand{\EqSrc}{\ensuremath{\mathrm{EqSrc}}}
\newtheorem{definition}{Definition}
\newtheorem{proposition}{Proposition}

\title{P3-T02: A Typed Source-Extension Declaration Lattice\\
for the \EqSrc{} Selector Obstruction}
\author{Candidate Constructor packet RT-20260720-023}
\date{2026-07-21}

\begin{document}
\maketitle

\begin{center}
\fbox{\begin{minipage}{0.93\linewidth}
\textbf{Artifact identity:} EQSRC-SOURCE-EXTENSION-LATTICE-V1.
\textbf{Status:} draft/control, proposal-only, source-extension data.
This packet constructs and tests candidate extension records.  It does not
adopt an extension, edit canonical ontology, establish physical gauge or
admissibility, discharge general \(\EqSrc\), change the Distance-to-GR ledger,
promote a benchmark, or authorize publication.
\end{minipage}}
\end{center}

\section{Question, baseline, and result}

The fixed P2-T03 theorem reduces objectwise deterministic selector existence
on a finite structural-automorphism core to a fixed-locus problem.  P3-T01
then distinguishes added source structure from representative irrelevance,
dynamics-mediated selection, probability-valued selection, and scoped
obstruction.  P3-T02 asks a narrower constructive question:

\begin{quote}
Which typed assumptions can be declared as candidate source extensions, how
are their assumption burdens partially ordered, and what finite calculation
shows whether a declared extension removes, preserves, or changes the selector
obstruction?
\end{quote}

The result has three parts.  First, an eight-atom Boolean object gives an exact
syntactic declaration envelope.  Second, a well-typed, target-dependent reduct
predicate cuts out the scientifically meaningful extension subposet.  Third,
an extension selects a deterministic point on a finite object exactly when
its compatible stabilizer-fixed locus is a singleton.  None of these results
derives the declared extension from the current \AE ther-flow ontology.

\section{Typed declaration atoms}

Let
\[
 {\cal A}=\{R,O,Q,G,M,I,B,D\}.
\]
The codes are bookkeeping labels with the following exact meanings.

\begin{longtable}{>{\raggedright\arraybackslash}p{0.07\linewidth}
                  >{\raggedright\arraybackslash}p{0.31\linewidth}
                  >{\raggedright\arraybackslash}p{0.51\linewidth}}
\toprule
Code & Declared datum or law & Boundary that remains open\\
\midrule
\endhead
\(R\) & distinguished root, ordered block, or partition provenance
& source derivation, preserving morphisms, and physical provenance\\
\(O\) & orientation, signed mark, associated line, or fibre datum
& an orientation-to-line map where needed, physical handedness, and gauge status\\
\(Q\) & differential, quotient discriminator, or discriminating readout law
& derivation of the differential or discriminator and its variation law\\
\(G\) & absolute grade, rank, order label, or declared minimum
& derivation of the grading; a parity grade need not have a unique minimum\\
\(M\) & normalized measure, stochastic kernel, or readout law
& deterministic point selection, operational meaning, and calibration\\
\(I\) & initial state, preparation, or history sector
& a universal law or a source-derived preparation mechanism\\
\(B\) & boundary value, asymptotic condition, or terminal condition
& boundary provenance, well-posedness, and universal applicability\\
\(D\) & transition relation, evolution law, ordered action, or flow
& symmetry breaking by itself, physical time, and empirical adequacy\\
\bottomrule
\end{longtable}

The separation \(Q\neq G\) is deliberate.  A chain differential or quotient
discriminator is not merely a grading, and an ordered monoid action is not
thereby physical time evolution.  Likewise, the three-line finite witness
below uses an explicit line mark.  An orientation-to-line claim would require
an additional equivariant map and is not inferred by terminology.

\section{The declaration envelope and admissible reduct order}

\begin{definition}[Syntactic declaration envelope]
The declaration envelope is
\[
 {\cal L}_{\mathrm{decl}}=(\Pow({\cal A}),\subseteq).
\]
For \(E,F\subseteq{\cal A}\),
\[
 E\wedge F=E\cap F,\qquad E\vee F=E\cup F.
\]
The rank of \(E\) is \(|E|\).  Rank is a count of declared atom types, not a
scientific cost, plausibility score, or adoption rank.
\end{definition}

\begin{proposition}[Exact finite envelope]
\({\cal L}_{\mathrm{decl}}\) is the Boolean lattice \(B_8\).  It has 256 nodes,
1024 directed cover edges, and rank sizes
\[
 (1,8,28,56,70,56,28,8,1).
\]
\end{proposition}

\begin{proof}
There are \(2^8=256\) subsets.  A directed cover adds exactly one absent atom.
Counting pairs \((E,a)\) with \(a\notin E\) gives
\(\sum_E(8-|E|)=8\,2^7=1024\).  Rank \(k\) has
\(\binom{8}{k}\) elements, giving the displayed sequence.  Intersection,
union, and complement in \({\cal A}\) supply the Boolean operations.
\end{proof}

\begin{definition}[Well-typed admissible subposet]
Fix a source baseline \(X\) and requested target type \(T\).  Let
\[
 {\cal L}_{\mathrm{adm}}(X,T)
 =\{E\in{\cal L}_{\mathrm{decl}}:\operatorname{WellTyped}_{X,T}(E)\}.
\]
For realized extensions \(e_E,e_F\), scientific assumption strength requires
a target-neutral reduct \(e_F\to e_E\) over the same baseline and target.
Boolean inclusion is a necessary declaration check, not by itself proof of
such a reduct.
\end{definition}

Measures require measurable structure; states and boundaries require their
carriers; dynamics requires a state space and transition semantics; and
orientation-to-line claims require an explicit map.  Consequently
\({\cal L}_{\mathrm{adm}}(X,T)\) need not be all of \(B_8\) and need not be
closed under Boolean meet or join.  The Boolean theorem is therefore an exact
accounting result, not a physical-realizability theorem.

\section{Compatible fixed-locus criterion}

Let a finite structural automorphism group \(K_X\) act on an eligible
deterministic choice set \({\cal D}(X)\).  For a realized extension package
\(e_E=(e_a)_{a\in E}\), define
\[
 H_E=\Stab_{K_X}(e_E)
     =\bigcap_{a\in E}\Stab_{K_X}(e_a).
\]
Let \(C_E(e_E)\subseteq{\cal D}(X)\) be the choices compatible with the
declared extension.  Compatibility must be \(H_E\)-stable.  Define
\[
 \Sigma(E,e_E)=C_E(e_E)\cap\Fix_{{\cal D}(X)}(H_E).
\]

\begin{proposition}[Finite compatible-selector trichotomy]
On the one-object finite groupoid of the realized extended object,
\[
\begin{array}{rcl}
|\Sigma(E,e_E)|=0 &\Longleftrightarrow& \text{no compatible deterministic selector},\\
|\Sigma(E,e_E)|=1 &\Longleftrightarrow& \text{one unique compatible deterministic selector},\\
|\Sigma(E,e_E)|>1 &\Longleftrightarrow& \text{residual deterministic underdetermination}.
\end{array}
\]
\end{proposition}

\begin{proof}
A deterministic choice on the extended object is natural under its structural
automorphisms exactly when it is fixed by every element of \(H_E\).  It is
admissible for the declared extension exactly when it lies in \(C_E(e_E)\).
Thus the set of compatible natural choices is precisely \(\Sigma(E,e_E)\);
the three conclusions follow from its cardinality.
\end{proof}

\paragraph{Essential non-monotonicity guard.}
If \(E\subseteq F\), then \(H_F\leq H_E\), hence
\(\Fix(H_E)\subseteq\Fix(H_F)\): reducing symmetry can enlarge the raw fixed
locus.  Uniqueness is therefore not monotone in declaration inclusion.
Compatibility must perform substantive work.  For noninvertible morphisms the
full P2-T03 equalizer guard remains necessary; the stabilizer formula here is
only the finite group-action case.

\paragraph{Different target types.}
An invariant probability law can be unique while the deterministic point
fixed locus is empty.  A relation image can likewise be unique while its point
representatives are not.  Point, probability, and relation targets must not be
interchanged.  P3-T03 owns representative-irrelevance proofs, and P3-T04 owns
dynamical and probabilistic selection theorems; this packet executes neither.

\section{Multidimensional assumption costs}

Each atom has a ten-component burden vector recorded in
\texttt{eqsrc\_source\_extension\_assumption\_costs\_v1.csv}.  Components
separately track source-signature, selector/readout, state/boundary, dynamics,
probability, variation-class, functoriality, empirical-calibration,
target-import, and protected-authority burdens.  Levels \(0,1,2,3\) are
within-axis declarations only.  Package aggregation is componentwise maximum.
No component sum, weighted score, or lexicographic preference is authorized.

Thus, for example,
\[
 \{Q\}\parallel\{G\},\qquad
 \{M\}\parallel\{D\},\qquad
 \{I\}\parallel\{B\},\qquad
 \{R,O\}\parallel\{R,D\}.
\]
These are genuine declaration incomparabilities.  They are not claims that
the packages are equally plausible or empirically supported.

\section{Finite witnesses}

The machine witness artifact reproduces the following exact values.

\begin{longtable}{>{\raggedright\arraybackslash}p{0.23\linewidth}
                  >{\raggedright\arraybackslash}p{0.11\linewidth}
                  >{\raggedright\arraybackslash}p{0.13\linewidth}
                  >{\raggedright\arraybackslash}p{0.13\linewidth}
                  >{\raggedright\arraybackslash}p{0.29\linewidth}}
\toprule
Witness & Package & \(|K_X|\) & \(|H_E|\) & Result\\
\midrule
\endhead
\texttt{W1} \(C_2\) root swap & \(R\) & 2 & 1
& one compatible fixed point; root provenance remains absent\\
\texttt{W2} three explicit lines & \(O\) & 6 & 2
& the marked line is the unique compatible fixed line\\
\texttt{W3} \(C_4\) absolute grading & \(G\) & 4 & 1
& unique declared grade-zero point; absolute zero is supplied\\
\texttt{W4} \(C_3\) uniform measure & \(M\) & 3 & 3
& no fixed point, but one invariant uniform distribution\\
\texttt{W5} symmetric two-branch dynamics & \(D\) & 2 & 2
& no fixed point; dynamics alone does not select a branch\\
\texttt{W6} state plus dynamics & \(I,D\) & 2 & 1
& one linked branch on the finite model\\
\texttt{W7} arbitrary \(C_3\) mark & \(R\) & 3 & 1
& formal singleton, rejected as a provenance/admissibility control\\
\texttt{W8} \(C_2\) discriminator fibre & \(Q\) & 2 & 1
& one declared zero-fibre point; discriminator provenance remains absent\\
\bottomrule
\end{longtable}

For \texttt{W2}, \(K_X\cong\mathrm{GL}(2,\mathbb F_2)\cong S_3\) acts
transitively on the three one-dimensional subspaces of \(\mathbb F_2^2\).
Orbit--stabilizer gives a marked-line stabilizer of order \(6/3=2\), and that
stabilizer fixes exactly the marked line.  This is a line-mark calculation,
not a physical-orientation claim.

For \texttt{W4}, \(C_3\) acts regularly on three points, so no point is fixed.
The normalized invariant weights satisfy
\[
 (w_0,w_1,w_2)=(w_2,w_0,w_1),\qquad
 w_0+w_1+w_2=1,
\]
and hence \(w_0=w_1=w_2=1/3\).  This proves one invariant distribution only.

\section{Placement of the five frozen historical families}

The table gives candidate cores minimal only on the named task-local finite
witness and within the eight-atom library.  The full family obligations in
\texttt{eqsrc\_source\_extension\_historical\_placements\_v1.yaml} remain
unmet.  No row reopens a frozen family.

\begin{longtable}{>{\raggedright\arraybackslash}p{0.33\linewidth}
                  >{\raggedright\arraybackslash}p{0.13\linewidth}
                  >{\raggedright\arraybackslash}p{0.43\linewidth}}
\toprule
Historical family & Candidate core & Principal obligations still absent\\
\midrule
\endhead
Intrinsic discriminator & \(\{Q\}\)
& derivation of differential/discriminator, morphisms, physical variations, and non-finite robustness\\
Cycle-boundary line & \(\{O\}\)
& derivation of oriented mark and mark-to-line law, variation scope, and physical meaning\\
Orientation torsor & \(\{R,O\}\)
& partition provenance, equivariant line selection, physical morphisms, and variations\\
Rooted partition & \(\{R,D\}\)
& transition projection, root provenance, path-parity admission, functoriality, and variations\\
Graded orbit root & \(\{G,D\}\)
& ordered action, full component, positivity, unique minimum, common grading, morphisms, and variations\\
\bottomrule
\end{longtable}

The last row does not double-count \(R\): if the declared ordered action and
grading actually prove a unique global minimum, the root is derived
conditionally inside that package.  The proof burden remains open.

\section{Distance-to-GR and freeze evaluation}

\begin{longtable}{>{\raggedright\arraybackslash}p{0.28\linewidth}
                  >{\raggedright\arraybackslash}p{0.62\linewidth}}
\toprule
Required field & P3-T02 result\\
\midrule
\endhead
Target derivation milestone & \texttt{source\_equivalence\_eqsrc}\\
Milestone burden & make the mathematical and ontological price of each selector repair explicit\\
New mathematical payload & exact \(B_8\) declaration envelope; well-typed reduct guard; compatible fixed-locus trichotomy; eight finite calculations\\
Existence status & task-local declaration object and finite witnesses exist\\
Uniqueness status & unique compatible selectors exist only on specified witnesses; no unique scientific extension is selected\\
Naturality/covariance status & objectwise finite automorphism checks only; full naturality, noninvertible equalizers, and physical covariance remain open\\
Source derivation status & current ontology does not derive \(R,O,Q,G,M,I,B,\) or \(D\) as adopted structures\\
Bridge attempt status & candidate extension space constructed; no extension adopted and no general \(\EqSrc\) bridge completed\\
Obstruction location & source-side provenance, typing, admissible morphisms/variations, and protected adoption\\
Distance-to-GR delta & none; the scientific ledger remains unchanged\\
\bottomrule
\end{longtable}

The repeated-family freeze remains in force.  P3-T02 is materially distinct
because it constructs a typed cross-family accounting object and finite
criterion rather than another family candidate.  Its result is
\texttt{blocked\_adoption\_open\_continuation}: adoption is human-gated, while
same-milestone theorem work remains open.  The packet proves neither that an
extension is impossible nor that any extension is physically preferred.

\section{Conclusion}

The smallest defensible conclusion is:

\begin{quote}
Under a declared, well-typed finite extension and compatible-choice functor,
the realized object has one deterministic selector exactly when its compatible
stabilizer-fixed locus is a singleton.  The eight-atom Boolean object
exhaustively records declaration combinations, while the admissible reduct
subposet and multidimensional cost vectors preserve typing and genuine
incomparability.  Current ontology does not derive or adopt any of these
extensions.
\end{quote}

P3-T03 is the next dependency-ready theorem packet.  It may ask when different
representatives become irrelevant at a declared target, but it must not infer
physical gauge or observable equivalence from structural symmetry alone.

\end{document}
